Kubernetes ist eine Plattform zur Orchestrierung von Containern und wird 
hauptsächlich für den Betrieb verteilter, skalierbarer Anwendungen eingesetzt. 
Im Unterschied zu Docker Compose, das mehrere Container auf einem einzelnen 
Host verwaltet, ist Kubernetes für den Betrieb in Cluster-Umgebungen mit 
mehreren Knoten ausgelegt \cite{kubernetes_docs}.

\subsubsection{Wie Kubernetes funktioniert}

Kubernetes abstrahiert Container in sogenannte \textbf{Pods}, die einen 
oder mehrere eng gekoppelte Container umfassen können. Mehrere Pods werden 
über \textbf{Services} zugänglich gemacht, während \textbf{Deployments} 
und \textbf{ReplicaSets} dafür sorgen, dass immer die gewünschte Anzahl 
an laufenden Instanzen vorhanden ist \cite{kubernetes_docs}.

Ein zentrales Merkmal von Kubernetes ist die deklarative Konfiguration: 
Man definiert den gewünschten Zustand, zum Beispiel dass immer fünf 
Instanzen eines Services laufen sollen, und Kubernetes sorgt selbstständig 
dafür, diesen Zustand herzustellen und aufrechtzuerhalten. Fällt eine 
Instanz aus, startet Kubernetes automatisch eine neue \cite{kubernetes_docs}.

\subsubsection{Ab wann macht Kubernetes Sinn}

Der eigentliche Vorteil von Kubernetes kommt erst bei größeren Systemen 
zum Tragen. Ein konkretes Beispiel: Läuft eine Anwendung zunächst mit 
einem einzigen Nginx-Container und steigt die Last plötzlich stark an, 
kann Kubernetes automatisch auf zwanzig Instanzen hochskalieren und die 
Last gleichmäßig verteilen. Sinkt die Last wieder, werden die überschüssigen 
Instanzen automatisch beendet \cite{kubernetes_docs}. Mit Docker Compose 
auf einer einzelnen virtuellen Maschine ist das nicht möglich.

Kubernetes eignet sich daher besonders für:
\begin{itemize}
    \item Anwendungen mit stark schwankender Last
    \item Microservice-Architekturen mit vielen unabhängigen Services
    \item Systeme mit hohen Anforderungen an Verfügbarkeit und 
    Ausfallsicherheit
    \item Cloud-Umgebungen mit mehreren Servern oder Rechenzentren
\end{itemize}

\subsubsection{Warum wir Kubernetes nicht verwenden}

Für LearnHub wurde bewusst auf Kubernetes verzichtet. Die Anwendung 
besteht aus einer überschaubaren Anzahl an Services und wird auf einer 
einzelnen virtuellen Maschine betrieben. In diesem Umfang bringt 
Kubernetes keinen Mehrwert, sondern nur zusätzliche Komplexität: Ein 
Kubernetes-Cluster benötigt eine Control Plane, Worker Nodes und ein 
eigenes Netzwerkmodell, was den Konfigurationsaufwand deutlich erhöht 
\cite{kubernetes_docs}.

Docker Compose ist für diesen Anwendungsfall die einfachere und 
wartungsfreundlichere Lösung. Sollte LearnHub in Zukunft stark wachsen 
und auf mehreren Servern betrieben werden müssen, wäre eine Migration 
zu Kubernetes möglich, da die Services bereits klar voneinander getrennt 
und vollständig containerisiert sind.